\chapter{Минимальные данные центрального гамильтониана дополнительных масс}
\label{ch:v8-extra-edge-mass-minimal-central-hamiltonian-data}

\section{Факторизация трёхуровневого центра}

На Real-завершённых блоках рангов $(4,6,6)$ общий сохраняющий gauge-типы
центральный гамильтониан имеет вид
\begin{equation}
 h=\varepsilon_YP_Y+\varepsilon_uP_u+\varepsilon_dP_d.
 \label{eq:v8-extra-mass-minimal-hamiltonian}
\end{equation}
Общий сдвиг энергии не меняет состояние Гиббса:
\begin{equation}
 \rho_{h+bI}=\frac{e^{-\beta(h+bI)}}{\Tr e^{-\beta(h+bI)}}
 =\frac{e^{-\beta h}}{\Tr e^{-\beta h}}=\rho_h.
 \label{eq:v8-extra-mass-hamiltonian-shift}
\end{equation}
Поэтому после quotient по $I$ остаются ровно две щели
\begin{equation}
 \Delta_u=\varepsilon_u-\varepsilon_Y,
 \qquad
 \Delta_d=\varepsilon_d-\varepsilon_Y,
 \label{eq:v8-extra-mass-two-gaps}
\end{equation}
или две безразмерные координаты
\begin{equation}
 \theta_u=\beta\Delta_u,
 \qquad
 \theta_d=\beta\Delta_d.
 \label{eq:v8-extra-mass-two-dimensionless-gaps}
\end{equation}

\section{Точная карта на внутренность симплекса}

Проективные центральные веса равны
\begin{equation}
 (p_Y,p_u,p_d)=\frac{(1,e^{-\theta_u},e^{-\theta_d})}
 {1+e^{-\theta_u}+e^{-\theta_d}}.
 \label{eq:v8-extra-mass-softmax-map}
\end{equation}
Обратная карта имеет вид
\begin{equation}
 \theta_u=\log\frac{p_Y}{p_u},
 \qquad
 \theta_d=\log\frac{p_Y}{p_d}.
 \label{eq:v8-extra-mass-softmax-inverse}
\end{equation}
Её якобиан невырожден во всей внутренности:
\begin{equation}
 \det\frac{\partial(p_u,p_d)}{\partial(\theta_u,\theta_d)}
 =p_Yp_up_d>0,
 \qquad
 \det J\big|_{\theta=0}=\frac1{27}.
 \label{eq:v8-extra-mass-softmax-jacobian}
\end{equation}
Следовательно, это диффеоморфизм $\mathbb R^2\to\operatorname{int}\Delta^2$:
две координаты одновременно необходимы и достаточны.

Три точных ориентира равны
\begin{equation}
 \begin{array}{c|c}
 (\theta_u,\theta_d)&(p_Y,p_u,p_d)\\ \hline
 (0,0)&(1/3,1/3,1/3)\\
 (\log(3/2),\log(3/2))&(3/7,2/7,2/7)\\
 (\log(3/2),\log3)&(1/2,1/3,1/6)
 \end{array}
 \label{eq:v8-extra-mass-gap-witnesses}
\end{equation}
Первая точка есть счётный trace, вторая --- равные edge-массы, третья
доказывает необходимость двух независимых, а не одной общей щели.

\section{Gibbs-вариация и минимальность}

В координатах $p_Y=1-p_u-p_d$ свободная энергия равна
\begin{equation}
 \Phi_{\theta}(p)=\theta_up_u+\theta_dp_d
 +p_Y\log p_Y+p_u\log p_u+p_d\log p_d.
 \label{eq:v8-extra-mass-free-energy}
\end{equation}
Её гессиан положительно определён:
\begin{equation}
 \det\Hess\Phi_{\theta}=\frac1{p_Yp_up_d}>0,
 \qquad
 \Hess\Phi_0\big|_{p=(1/3,1/3,1/3)}=
 \begin{pmatrix}6&3\\3&6\end{pmatrix}.
 \label{eq:v8-extra-mass-free-energy-hessian}
\end{equation}
Поэтому каждая пара щелей имеет единственный Gibbs-минимум. Это доказывает
достаточность параметров, но не их parent-origin; detailed balance также
формулируется относительно уже выбранного верного состояния
\cite{v8-extra-mass-fagnola}.

Бесследовый представитель
\begin{equation}
 4\varepsilon_Y+6\varepsilon_u+6\varepsilon_d=0
 \label{eq:v8-extra-mass-traceless-representative}
\end{equation}
лишь фиксирует общий сдвиг и оставляет двумерное ядро. Он не выбирает
$\Delta_u,\Delta_d$.

\section{Неидентифицируемые масштабы}

Равновесное состояние инвариантно относительно общей масштабной орбиты
\begin{equation}
 (\beta,\Delta_u,\Delta_d)\longmapsto
 (\beta/c,c\Delta_u,c\Delta_d),\qquad c>0.
 \label{eq:v8-extra-mass-thermal-scale-orbit}
\end{equation}
Поэтому веса определяют только $(\theta_u,\theta_d)$, но не температуру и
две физические энергии по отдельности.

Сами коэффициенты стабилизатора параметризуются как
\begin{equation}
 \boldsymbol\mu=\alpha(2p_Y,3p_u,3p_d),
 \qquad \alpha>0.
 \label{eq:v8-extra-mass-scale-parameterization}
\end{equation}
Пара щелей выбирает отношения, но не общий масштаб $\alpha$. Временная
скорость релаксации дополнительно требует спектральной плотности среды, что
не следует из Gibbs-состояния \cite{v8-extra-mass-davies}.

Итоговые реестры равны
\begin{equation}
 \begin{aligned}
 N_{\rm architecture}&=7/7,
 &N_{\rm relative\ gap\ origin}&=0/2,\\
 N_{\rm energy\ scale}&=0/1,
 &N_{\rm mass\ scale}=N_{\rm relaxation\ scale}&=0/1.
 \end{aligned}
 \label{eq:v8-extra-mass-minimal-data-ledgers}
\end{equation}

\begin{theorem}[Минимальность пары центральных щелей]
С точностью до общего сдвига всякий type-сохраняющий центральный гамильтониан
трёх дополнительных edge-блоков задаётся двумя энергетическими разностями.
Пара $(\theta_u,\theta_d)=(\beta\Delta_u,\beta\Delta_d)$ диффеоморфно
параметризует всю внутренность положительного trace-симплекса и является
минимальным достаточным набором относительных данных.
\end{theorem}

\begin{nogocriterion}[Равновесные веса не задают физические масштабы]
Нельзя извлечь из $(p_Y,p_u,p_d)$ раздельные $\beta$, $\Delta_u$ и
$\Delta_d$, поскольку они лежат на общей масштабной орбите. Нельзя также
отождествлять безразмерные щели с общим масштабом масс $\alpha$ или временем
релаксации.
\end{nogocriterion}

\begin{protocol}\begin{enumerate}
\item центральных энергетических уровней: \textbf{три};
\item физически незначимых общих сдвигов: \textbf{один};
\item независимых щелей: \textbf{две};
\item минимальные координаты: \textbf{$(\theta_u,\theta_d)$};
\item карта на $\operatorname{int}\Delta^2$: \textbf{диффеоморфизм};
\item якобиан: \textbf{$p_Yp_up_d>0$};
\item Gibbs-минимум: \textbf{единственный};
\item архитектура: \textbf{$7/7$};
\item parent-origin двух щелей: \textbf{$0/2$};
\item абсолютная энергия: \textbf{$0/1$};
\item общий масштаб масс: \textbf{$0/1$};
\item скорость релаксации: \textbf{$0/1$};
\item следующий гейт: \textbf{parent-origin центрального гамильтониана}.
\end{enumerate}\end{protocol}

Машинный сертификат сохранён в
\path{s2t_v8_baryon_c0_singlet_triplet_central_gap_isotypic_channel_extra_edge_mass_minimal_central_hamiltonian_data_gate_results.json}.

\begin{thebibliography}{9}
\bibitem{v8-extra-mass-fagnola}
F.~Fagnola, V.~Umanit\`a,
\emph{Generators of Detailed Balance Quantum Markov Semigroups},
arXiv:0707.2147.
\bibitem{v8-extra-mass-davies}
E.~B.~Davies,
\emph{Markovian Master Equations}, Commun. Math. Phys. 39 (1974), 91--110.
\bibitem{v8-extra-mass-gibbs}
H.~Moriya,
\emph{Gibbs Variational Formula for Thermal Equilibrium States in Terms of
Quantum Relative Entropy Density}, arXiv:2002.04253.
\end{thebibliography}